\clearpage
\thispagestyle{empty}
\pdfbookmark[0]{Ein neues Hilfsmittel für die magnetischen Beobachtungen. Ausgabe und Quellen}{anzeige-1837-10-30-sources}
\begin{center}\Large\textbf{Ein neues Hilfsmittel für die magnetischen Beobachtungen}\end{center}

Diese Anzeige zu Carl Friedrich Gauss' Vortrag vom 19. September 1837 erschien am 30. Oktober 1837 in den \textit{Göttingischen gelehrten Anzeigen}. Die KI-Übertragung in modernes Deutsch mit KI-erstelltem Satz folgt der gesammelten Ausgabe in den \textit{Werken}, Band V (Göttingen, 1867), S.~352--356. Die Längenangaben und das Gewicht des Magneten stehen in den historischen Einheiten Meile, Fuß und Pfund.

\medskip\noindent
Gauß' eigenständige Abhandlung über dasselbe Instrument erschien in den \textit{Resultaten} für 1837, die 1838 veröffentlicht wurden. William Francis übersetzte diese Abhandlung für Richard Taylors \textit{Scientific Memoirs}, Bd.~II (1841), S.~252--267. Major Sabine und Professor Lloyd revidierten die Übersetzung. Deren Anfang wurde in André Koch Torres Assis (Hrsg.), \textit{Wilhelm Weber's Main Works on Electrodynamics Translated into English}, Bd.~I (2021), Kap.~11, zum Vergleich herangezogen.

\medskip\noindent
\href{https://gdz.sub.uni-goettingen.de/id/PPN236006339}{Gesammelte Ausgabe (1867)}\quad
\href{https://www.ifi.unicamp.br/~assis/Weber-in-English-Vol-1.pdf\#page=126}{Eigenständige Abhandlung auf Englisch}
\clearpage
